% Часть 4: соглашения о нормировке, сводка, итог.

\subsection{Соглашения о нормировке и сверка всех коэффициентов}
\label{sec:kd_conv}

\paragraph{Зачем.} Численные множители вроде $\tfrac\pi2$, $\tfrac\pi4$, $\tfrac14$, $\tfrac3{16}$ в
разных статьях о Кондо-обмене часто различаются не из-за ошибок, а из-за
соглашений: как записан гамильтониан обмена и что называют плотностью
состояний. Перечислим соглашения, найдём, какую нормировку подразумевает
каждая формула обеих статей, и отделим соглашения от настоящих
несогласованностей.

\paragraph{Какие бывают соглашения.}
\begin{enumerate}
\item \textbf{Спин электрона: $J\,\mathbf S\cdot\mathbf s$ или $J'\,\mathbf S\cdot\boldsymbol\sigma$.} $J'=J/2$, поэтому
  коэффициенты при $(\rho J')^2$ в 4 раза больше, чем при $(\rho J)^2$.
\item \textbf{Плотность состояний на одну проекцию спина ($\rho_\sigma$) или на обе
  ($\rho=2\rho_\sigma$).} $(\rho J)^2=4(\rho_\sigma J)^2$: коэффициенты при $(\rho J)^2$ в 4 раза
  меньше.
\item \textbf{Скорость или полуширина.} Скорость ухода населённости $\Gamma$ и
  полуширина линии $\Gamma/2$ отличаются вдвое. Это касается только скоростей, не
  сдвигов.
\item \textbf{Единицы} $\hbar=1$, $k_B=1$ на безразмерные коэффициенты не влияют.
\end{enumerate}
Знак $J$ (ферро- или антиферромагнитный обмен) во второй порядок не входит:
везде стоит $J^2$.

\paragraph{Что фиксировано в наших статьях.} Обе статьи пишут обмен через
$\mathbf s=\boldsymbol\sigma/2$ (Surf.\ Sci.\ ур.~3, PRB ур.~4--5, раздел~\ref{sec:kd_notation}), так что
соглашение~1 у них одно. Свободным остаётся соглашение о плотности
состояний:
\begin{paperbox}{PRB, перед ур.~19; Surf.\ Sci., перед ур.~14}
\[
\rho(\epsilon)=\sum_{\mathbf k,\sigma}\delta(\epsilon_{\mathbf k\sigma}-\epsilon)\quad\text{(PRB: сумма по обеим проекциям)};
\qquad \rho=\frac1W\quad\text{(Surf.\ Sci.)}.
\]
\end{paperbox}

\paragraph{Как читать таблицу.} Для каждой формулы запишем коэффициент,
который получается из~\eqref{eq:kd_rate}, \eqref{eq:kd_shift}, \eqref{eq:kd_shift34}
и~\eqref{eq:kd_ad} при каждом соглашении о $\rho$, и сравним с напечатанным. Мера ---
\begin{resbox}[Мера нормировки формулы]
\[
x=\frac{\text{правильный коэффициент при }(\rho_\sigma J)^2}{\text{напечатанный коэффициент при }(\rho J)^2},
\qquad (\rho J)^2_{\text{статьи}}=x\,(\rho_\sigma J)^2 .
\]
$\rho$ на одну проекцию даёт $x=1$, на обе --- $x=4$. Формулы в одной нормировке
обязаны иметь одинаковое $x$.
\end{resbox}

Дефазировка в Surf.\ Sci.\ (нужна для строки таблицы):
\begin{paperbox}{Surf.\ Sci., ур.~16}
\[
\gamma^{\rm ad}_{nm}=\frac{\pi}{4\hbar}(\rho J)^2k_BT\sum_{\sigma\sigma'}\sum_{a=x,y,z}
\Big|\tau^a_{\sigma\sigma'}\langle n|S_a|n\rangle-\tau^a_{\sigma'\sigma}\langle m|S_a|m\rangle\Big|^2 .
\]
\end{paperbox}
Для состояний $\pm S$ работает только $a=z$:
$\sum_{\sigma\sigma'}|\tau^z_{\sigma\sigma'}(S-(-S))|^2=2\cdot4S^2$, и
$\tfrac\pi4\cdot8S^2=2\pi S^2$.

\begin{center}\small
\begin{tabular}{p{4.7cm}cccc}
\toprule
Формула & напечатано & при $\rho_\sigma$ & при $\rho=2\rho_\sigma$ & $x$ \\
\midrule
SS ур.~14, скорость & $\tfrac\pi2$ & $\pi$ & $\tfrac\pi4$ & 2 \\
SS ур.~20--22, сдвиг (поле) & $\tfrac14$ & $\tfrac12$ & $\tfrac18$ & 2 \\
SS ур.~23, сдвиг $S=\tfrac32$ & $\tfrac3{16}$ & $\tfrac32$ & $\tfrac38$ & 8 \\
SS ур.~24, сдвиг $S=1$ & $\tfrac32$ & $\tfrac32$ & $\tfrac38$ & 1 \\
SS ур.~16, дефазир.\ $\pm S$ & $2\pi k_BTS^2$ & $2\pi k_BTS^2$ & $\tfrac\pi2k_BTS^2$ & 1 \\
PRB ур.~20, 22, скорость & $\tfrac\pi4$, $\tfrac\pi8$ & $\pi$, $\tfrac\pi2$ & $\tfrac\pi4$, $\tfrac\pi8$ & 4 \\
PRB ур.~38--39, дефазир. & $\tfrac\pi8$, $\tfrac\pi2k_BTS^2$ & $\tfrac\pi2$, $2\pi k_BTS^2$ & $\tfrac\pi8$, $\tfrac\pi2k_BTS^2$ & 4 \\
\bottomrule
\end{tabular}
\end{center}
Откуда столбцы: скорость $\pi(\rho_\sigma J)^2=\tfrac\pi4(\rho J)^2$ (раздел~\ref{sec:kd_rate}); сдвиг в
поле $\tfrac12(\rho_\sigma J)^2=\tfrac18(\rho J)^2$, анизотропия --- втрое больше
(разделы~\ref{sec:kd_shift}, \ref{sec:kd_Bn}); дефазировка $\pm S$ ---
$\tfrac{\pi(\rho J)^2}{8}k_BT\cdot4S^2$ (раздел~\ref{sec:kd_deph}).

\begin{paperbox}{PRB, ур.~26 (одномерная зона)}
\[
\Gamma_{MM'}=\frac{\pi(\rho J)^2}{32\hbar}\,G(\Delta)\,2\big[\chi_{M,M'}(0)+\chi_{M,M'}(2k_Fd)\big].
\]
\end{paperbox}
\begin{cmpbox}
При $d=0$ для спина $\tfrac12$ ($\chi(0)=1$): $\tfrac{\pi}{32}\cdot2\cdot2=\tfrac\pi8$, то есть $\Gamma_0$ ур.~22.
Согласовано.
\end{cmpbox}

\begin{cmpbox}[Что следует из таблицы]
\begin{itemize}
\item \textbf{PRB:} все формулы имеют $x=4$ --- плотность на обе проекции, как и
  определено в статье. Нормировка выдержана; в ней написан наш код.
\item \textbf{Surf.\ Sci., ур.~14, 20, 21, 22:} одинаковое $x=2$, формулы согласованы
  между собой (то же даёт отношение $\tfrac1{2\pi}$, раздел~\ref{sec:kd_shift}). Само
  $x=2$ не совпадает ни с одним соглашением: при гамильтониане ур.~3 это
  половина результата для $\rho_\sigma$ или удвоенный для $\rho=2\rho_\sigma$. Это общий
  масштаб, на сравнения он не влияет. Прямая подстановка ур.~13 в ур.~12 с
  $\mathrm{Im}\,\tfrac1{x+i0}=-\pi\delta(x)$ даёт $2\pi(\rho J)^2$: общие ур.~12--13 записаны с
  точностью до числового множителя.
\item \textbf{Surf.\ Sci., ур.~23:} $x=8$, в 4 раза больше, чем у ур.~20--22 --- ровно
  переход между $\rho_\sigma$ и $2\rho_\sigma$. Формула взята из Oberg et al.\ и,
  по-видимому, сохранила её нормировку.
\item \textbf{Surf.\ Sci., ур.~24:} $x=1$, вдвое меньше, чем у ур.~20--22. Соглашения~1
  и~2 дают только множитель 4, а соглашение~3 к сдвигам не относится:
  множитель 2 соглашениями не объясняется.
\item \textbf{Surf.\ Sci., ур.~16:} $x=1$, вдвое меньше, чем у ур.~14. Ниже показано,
  что это несогласованность скорости и дефазировки.
\end{itemize}
\end{cmpbox}

\paragraph{Проверки, не зависящие от нормировки.} Чтобы отличить соглашение
от ошибки, нужны отношения, в которых $\rho$ и $J$ сокращаются.
\begin{resbox}[Три отношения, обязательные при любой нормировке]
\begin{enumerate}
\item Сдвиг : скорость $=\dfrac1{2\pi}$ (одна собственная энергия,
  раздел~\ref{sec:kd_shift}).
\item Анизотропия : поле $=3$ (двойной коммутатор, раздел~\ref{sec:kd_Bn}).
\item $T_1=T_2$ для спина $\tfrac12$ при $B\to0$: в изотропной ванне без поля нет
  выделенного направления, поэтому продольная и поперечная релаксация
  идут с одной скоростью.
\end{enumerate}
\end{resbox}
Проверка 3 для PRB. Скорости (ур.~20, $G(0)=k_BT$):
\[
\Gamma_{\uparrow\to\downarrow}=\Gamma_{\downarrow\to\uparrow}=\tfrac\pi4(\rho J)^2k_BT\cdot\tfrac12=\tfrac\pi8(\rho J)^2k_BT,
\qquad \frac1{T_1}=\Gamma_{\uparrow\to\downarrow}+\Gamma_{\downarrow\to\uparrow}=\tfrac\pi4(\rho J)^2k_BT .
\]
Затухание когерентности --- неадиабатическая часть плюс адиабатическая
(ур.~38, $\Delta S^z=1$):
\begin{paperbox}{PRB, ур.~14}
\[
\gamma^{\rm nonad}_{M,M'}=\frac12\Big(\sum_{N\ne M}\Gamma_{M,N}+\sum_{N\ne M'}\Gamma_{M',N}\Big).
\]
\end{paperbox}
\[
\frac1{T_2}=\underbrace{\tfrac12(\Gamma_{\uparrow\to\downarrow}+\Gamma_{\downarrow\to\uparrow})}_{\frac\pi8(\rho J)^2k_BT}
+\underbrace{\tfrac\pi8(\rho J)^2k_BT}_{\text{ур.~38}}=\tfrac\pi4(\rho J)^2k_BT=\frac1{T_1}.
\]
Для Surf.\ Sci.\ в нормировке ур.~14:
\[
\frac1{T_1}=\tfrac\pi2(\rho J)^2k_BT,\qquad
\frac1{T_2}=\underbrace{\tfrac\pi4(\rho J)^2k_BT}_{\text{неадиаб.}}+\underbrace{\tfrac\pi4\cdot2\,(\rho J)^2k_BT}_{\text{ур.~16}}
=\tfrac{3\pi}4(\rho J)^2k_BT\ne\frac1{T_1}.
\]
\begin{cmpbox}[Результат трёх проверок]
\begin{center}\small
\begin{tabular}{p{3.4cm}p{5.4cm}p{3.6cm}}
\toprule
Проверка & Surf.\ Sci. & PRB \\
\midrule
сдвиг : скорость $=\tfrac1{2\pi}$ & ур.~14 и 20: выполнено & --- \\
анизотропия : поле $=3$ & ур.~23: $\tfrac34$; ур.~24: 6 --- нарушено & --- \\
$T_1=T_2$ при $B\to0$ & ур.~14 и 16: нарушено (множитель 2) & ур.~20 и 38: выполнено \\
\bottomrule
\end{tabular}
\end{center}
Проверка $T_1=T_2$ в PRB подтверждает и множитель $\tfrac12$ в~\eqref{eq:kd_ad}.
\end{cmpbox}

\begin{cmpbox}[Итог по коэффициентам]
\begin{itemize}
\item Префакторы $\tfrac\pi2$ (Surf.\ Sci.) и $\tfrac\pi4$ (PRB) и общий масштаб ур.~20--22 ---
  различие соглашений, на результаты не влияющее.
\item Коэффициенты ур.~23 и~24 Surf.\ Sci.\ по отдельности совпадают с правильным
  ответом в каких-то нормировках, но не в той, в которой написаны
  остальные формулы статьи (отношение 3:1 нарушено).
\item Ур.~16 Surf.\ Sci.\ несогласовано с ур.~14 (проверка $T_1=T_2$). Дефазировку мы
  моделируем по PRB, где все проверки выполнены.
\item Дополнительные материалы Oberg et al., где задана нормировка ур.~23, нам
  недоступны; вывод о ней сделан по множителю~4.
\end{itemize}
\end{cmpbox}

\subsection{Сводка и итог}
\label{sec:kd_summary}

\begin{center}\small
\begin{tabular}{p{3.4cm}p{3.4cm}p{7.2cm}}
\toprule
Коэффициент & Где встречается & Откуда берётся \\
\midrule
$J^2/2$ & квадрат матричного элемента & $(J/2)^2$ от $\mathbf s=\boldsymbol\sigma/2$; $\times2$ от
$\mathrm{Tr}\,\sigma^a\sigma^b=2\delta_{ab}$ (сумма по спинам электрона) \\
$\pi(\rho_\sigma J)^2=\tfrac\pi4(\rho J)^2$ & скорость (PRB, код) & $2\pi$ золотого правила
$\times\ J^2/2\times\rho_\sigma^2$; $\rho=2\rho_\sigma$ \\
$\tfrac\pi2(\rho J)^2$ & скорость (Surf.\ Sci.) & общий масштаб ур.~14, 20--22 ($x=2$, раздел~\ref{sec:kd_conv}) \\
$G=\dfrac{\Delta}{1-e^{-\beta\Delta}}$ & энергетический множитель &
$\int f(\epsilon)[1-f(\epsilon+\Delta)]d\epsilon=\Delta[1+n_B(\Delta)]$ \\
$K_d$: $\cos^2$, $J_0^2$, $\mathrm{sinc}^2$ & интерференция атомов & $|\langle e^{ik_F\hat{\mathbf k}\cdot\mathbf r}\rangle|^2$ по
поверхности Ферми \\
$\tfrac12$ в $\Gamma_0$ & спин $1/2$ & $|\langle\downarrow|S^x|\uparrow\rangle|^2+|\langle\downarrow|S^y|\uparrow\rangle|^2=\tfrac14+\tfrac14$ \\
$\tfrac1{2\pi}$ & сдвиг / скорость & $\mathrm{Re}\,\Sigma$ и $-2\,\mathrm{Im}\,\Sigma$ одной собственной энергии \\
1 при $L=\ln(W/k_BT)$ & сдвиг уровней & $\int\!\!\int(x+y)^{-2}dx\,dy$, обрезание на $k_BT$ \\
$\tfrac14$ и $\tfrac34$ & сдвиг: поле, анизотропия & $\tfrac14\times\tfrac{k(k+1)}2$; двойной
коммутатор, $k=1,2$ \\
$1-K$, $1+K$ & димер & элементы $S^z$ (в синглет) и $S^{x,y}$ (внутри триплета) \\
$\tfrac12$ в $\gamma^{\rm ad}$ & дефазировка & действие $L_0\rho L_0^\dagger-\tfrac12\{L_0^\dagger L_0,\rho\}$ на
$\rho_{MM'}$; проверка $T_1=T_2$ при $B\to0$ \\
$\times4$ в $(\rho J)^2$ & разные статьи & $\rho=2\rho_\sigma$ или $\mathbf S\cdot\boldsymbol\sigma$ вместо $\mathbf S\cdot\mathbf s$ (раздел~\ref{sec:kd_conv}) \\
\bottomrule
\end{tabular}
\end{center}

\begin{center}\small
\begin{tabular}{p{3.6cm}p{5.6cm}p{4.8cm}}
\toprule
Расхождение & Чем проверено & Где ошибка \\
\midrule
Префакторы $\tfrac\pi2$ и $\tfrac\pi4$ & вывод из золотого правила; отношение
$\tfrac1{2\pi}$ внутри Surf.\ Sci. & нет ошибки: разные нормировки \\
Ур.~23 Surf.\ Sci. & отношение 3:1 к ур.~22; явный подсчёт для $S=\tfrac32$;
источник --- Oberg et al.\ 2014, ур.~(2) & множитель 4 = смена соглашения о $\rho$;
не пересчитано к нормировке статьи \\
Ур.~24 Surf.\ Sci. & отношение 3:1 к ур.~22; явный подсчёт для $S=1$; текст
статьи после ур.~24 & множитель 2, соглашениями не объясняется \\
Ур.~16 Surf.\ Sci. & $T_1=T_2$ при $B\to0$ вместе с ур.~14 & множитель 2 между
ур.~14 и~16; в PRB согласовано \\
Собственные значения $S=1$ & диагонализация & опечатка в знаке; на выводы не
влияет \\
Рис.~2(b) PRB, 2D и 3D & предел независимых спинов; панель (c);
ур.~31--32 & в рисунке статьи (множитель $\tfrac12$) \\
Ур.~16, рис.~4, прил.~C PRB & вывод из~\eqref{eq:kd_V}; приложение B;
ур.~38--39; независимость от начала координат & в знаке показателя ур.~16 \\
\bottomrule
\end{tabular}
\end{center}

\begin{cmpbox}[Итог]
Во всех содержательных расхождениях --- коэффициенты ур.~23--24 Surf.\ Sci.,
множитель $\tfrac12$ на рис.~2(b) PRB и знак в ур.~16 PRB --- наш код совпадает с той
версией, которая следует из исходных выражений самих статей и проходит
проверки, не зависящие от нормировки: отношение 3:1, предел независимых
спинов, независимость от начала координат, $T_1=T_2$ при $B\to0$, согласие с
ур.~38--39. Разница префакторов $\tfrac\pi2$ и $\tfrac\pi4$ --- соглашение о нормировке.
Коэффициент ур.~23 Surf.\ Sci.\ взят из работы Oberg et al.\ и отличается от
согласованного с ур.~22 значения ровно множителем 4, который даёт смена
соглашения о плотности состояний; для множителя 2 в ур.~24 и для рисунков
такого объяснения нет.
\end{cmpbox}
